% =====================================================================
%  TYPE: Postoperative instructions -- trigeminal RF rhizotomy (Gasserian ganglion)
%  Sections: patient block, instructions, follow-up, emergency
%  box, when to call, medication page. Based on Xarokopou 2026-08.
%  Written for the RIGHT side, branches V2-V3. For the left side,
%  search for 'δεξ' and 'αριστ' and swap them.
%
%  Copy, rename to "YYYYMMDD - surname - topic.tex", fill in and
%  compile with XeLaTeX. Toggle alternatives by moving the "%".
%  Text in [brackets] is a placeholder: replace or delete it.
% =====================================================================
\documentclass[10pt,greek]{article}
\usepackage[useregional]{datetime2}

%%%%%%%%%%%%%%%%%%%%% Language
\providecommand\DocLang{gr}

%%%%%%%%%%%%%%%%%%%%% Clinic (letterhead logo, footer address)
\providecommand\Clinic{Peristeri}
% \providecommand\Clinic{Mandra}
% \providecommand\Clinic{Bioclinic}
% \providecommand\Clinic{Psychikou}
% \providecommand\Clinic{Metropolitan}
% \providecommand\Clinic{Dunant}
% \providecommand\Clinic{Euroclinic}

%%%%%%%%%%%%%%%%%%%%% Document code for the QR: paste the code issued
%%%%%%%%%%%%%%%%%%%%% by doc.erns.eu (empty = no QR code)
\providecommand\DocCode{}

%%%%%%%%%%%%%%%%%%%%% Patient
\providecommand\PatientPrefix{Ο}
% \providecommand\PatientPrefix{Η κα}
\providecommand\PatientName{Παπαδόπουλος Γιώργος}
\providecommand\PatientAge{45}
% \providecommand\PatientAge{}

%%%%%%%%%%%%%%%%%%%%% Date
\providecommand\DocDate{\DTMdate{2026-01-15}}
% \providecommand\DocDate{\DTMdate{2026-05-04}}

%%%%%%%%%%%%%%%%%%%%% Title
\providecommand\DocTitle{Μετεγχειρητικές Οδηγίες}
% \providecommand\DocTitle{Ενημερωτικό Σημείωμα για τον Ασθενή}

%%%%%%%%%%%%%%%%%%%%% Closing (empty = no signature block)
\providecommand\DocClosing{Ο ιατρός}
% \providecommand\DocClosing{}

%%%%%%%%%%%%%%%%%%%%% Stamp and signature image
\providecommand\DocStamp{no}
% \providecommand\DocStamp{yes}

\input{evrtemplate}

\begin{document}

\begin{PatientInfo}
  \Info{Διάγνωση}{Νευραλγία τριδύμου δεξιά (Τιμή 1 κλάδος), διάρκειας Τιμή 2.}
  \Info{Χειρουργείο}{(DD/MM/YYYY) Διαδερμική ριζοτομή του γασσερίου γαγγλίου με ραδιοσυχνότητες (RF) δεξιά -- συνδυασμός \textbf{παλμικής (pulsed)} και \textbf{θερμικής (thermal)} βλάβης, στους κλάδους \V{V2} και \V{V3}.}
%   \Info{Συνοδά}{Τιμή 4}
\end{PatientInfo}

Ακολουθούν οδηγίες για την ομαλή σας ανάρρωση. \textbf{Οι οδηγίες 1 και 2 είναι οι πλέον κρίσιμες} και αφορούν την ασφάλειά σας τις πρώτες εβδομάδες.

\section*{Κρίσιμες Οδηγίες}
\begin{enumerate}[label=\arabic*.]
\setlength{\parskip}{6pt}
  \item \textbf{Πρόληψη πτώσεων}

Τα φάρμακα που λαμβάνετε (Lyrica, Tegretol) προκαλούν ζάλη, αστάθεια και υπνηλία, ιδίως τις πρώτες ημέρες.
% Σε συνδυασμό με τη νευρολογική σας κατάσταση, ο κίνδυνος πτώσης είναι \textbf{αυξημένος}.
       \begin{itemize}
\item \textbf{Μην μετακινείστε ποτέ μόνος/μόνη σας} τις πρώτες ημέρες -- πάντα με συνοδό.
\item Σηκωθείτε \textbf{αργά} από την καθιστή ή ξαπλωτή θέση. Καθίστε πρώτα στην άκρη του κρεβατιού για 1--2 λεπτά πριν σταθείτε όρθιος/όρθια.
% \item Χρησιμοποιείτε το βοήθημα βάδισης (μπαστούνι / περιπατητήρα) που ήδη χρησιμοποιείτε, ακόμη κι αν αισθάνεστε καλά.
\item Απομακρύνετε χαλάκια, καλώδια και εμπόδια από τους χώρους που κινείστε. Φροντίστε για επαρκή φωτισμό, ιδίως τη νύχτα.
\item Φοράτε κλειστά, αντιολισθητικά υποδήματα -- ποτέ παντόφλες χωρίς πτέρνα.
\item Στο μπάνιο χρησιμοποιείτε αντιολισθητικό χαλάκι και, εφόσον υπάρχει, κάθισμα ντους.
\item \textbf{Μην οδηγείτε} και μην χειρίζεστε μηχανήματα.
    \end{itemize}

  \item \textbf{Προστασία της περιοχής με μειωμένη αίσθηση}

Μετά την επέμβαση, η δεξιά παρειά, τα ούλα, το χείλος και τμήμα της γλώσσας ενδέχεται να είναι \textbf{μουδιασμένα}. Αυτό είναι \emph{αναμενόμενο} και αποτελεί μέρος του θεραπευτικού αποτελέσματος. Ωστόσο σημαίνει ότι \textbf{δεν θα αντιληφθείτε τραυματισμό ή έγκαυμα} στην περιοχή:
       \begin{itemize}
\item \textbf{Ελέγχετε τη θερμοκρασία} των ροφημάτων και του φαγητού με τα χείλη της \emph{αριστερής} πλευράς ή με το χέρι, πριν τα βάλετε στο στόμα.
\item \textbf{Μασάτε από την αριστερή πλευρά.} Τρώτε αργά, σε μικρές μπουκιές, για να μην δαγκώσετε την παρειά ή τη γλώσσα σας.
\item Μετά από κάθε γεύμα, \textbf{ελέγχετε το εσωτερικό της δεξιάς παρειάς} με καθρέφτη για πληγές ή δήγματα.
\item Ξεπλένετε το στόμα μετά τα γεύματα -- υπολείμματα τροφής μπορεί να παραμείνουν χωρίς να το αντιληφθείτε.
\item Προσοχή στο ξύρισμα / στην περιποίηση προσώπου στη δεξιά πλευρά.
\item Επισκεφθείτε οδοντίατρο εάν παρατηρήσετε πληγή που δεν επουλώνεται, ή εάν καθυστερήσει η επανεκτίμηση.
    \end{itemize}
\end{enumerate}

\section*{Γενικές Οδηγίες}
\begin{enumerate}[label=\arabic*.]
\setlength{\parskip}{6pt}
  \item \textbf{Αναμενόμενες αισθητικές μεταβολές}
       \begin{itemize}
\item \textbf{Μούδιασμα} στη δεξιά παρειά, άνω/κάτω χείλος, ούλα και σιαγόνα. Αναμενόμενο. Υποχωρεί μερικώς μέσα σε μερικές εβδομάδες.
\item Αίσθημα «μυρμηγκιάσματος», καύσου ή ηλεκτρικού ρεύματος στην περιοχή τις πρώτες εβδομάδες. Υποχωρεί σταδιακά.
\item Πιθανή ήπια ασυμμετρία στην αίσθηση του προσώπου.
\item \textbf{Ενημερώστε τον ιατρό} εάν το μούδιασμα εξελιχθεί σε \emph{ενοχλητικό, επίμονο καυστικό άλγος} σε περιοχή που είναι ταυτόχρονα αναίσθητη.
    \end{itemize}
  \item \textbf{Μάσηση -- Κροταφογναθική άρθρωση}
       \begin{itemize}
\item Είναι πιθανή \textbf{αδυναμία των μασητήρων μυών δεξιά} μετά από βλάβη του 3ου κλάδου, με δυσκολία στη μάσηση και ελαφρά απόκλιση της σιαγόνας.
\item Τις πρώτες ημέρες προτιμήστε \textbf{μαλακή τροφή} (σούπες, πουρές, γιαούρτι, αυγά, ψάρι).
\item Αποφύγετε σκληρές, τραγανές ή κολλώδεις τροφές (καρύδια, παξιμάδια, ξηρούς καρπούς, σκληρή τσίχλα).
\item Ήπιες ασκήσεις ανοίγματος/κλεισίματος του στόματος αρκετές φορές την ημέρα, εάν υπάρχει δυσκαμψία (τρισμός).
\item Η μυϊκή ισχύς συνήθως βελτιώνεται σταδιακά μέσα σε εβδομάδες έως μήνες.
    \end{itemize}
  \item \textbf{Προστασία του οφθαλμού}
       \begin{itemize}
\item Η βλάβη αφορούσε τους κλάδους V\textsubscript{2} και V\textsubscript{3} -- ο οφθαλμικός κλάδος (V\textsubscript{1}) δεν στοχεύθηκε και το αντανακλαστικό του κερατοειδούς αναμένεται ακέραιο.
\item Ωστόσο, \textbf{εάν παρατηρήσετε μούδιασμα στο μέτωπο, στο άνω βλέφαρο ή στη μύτη}, ή \textbf{εάν το δεξί μάτι δεν αισθάνεται} τη σκόνη ή τον αέρα, ενημερώστε \textbf{άμεσα} τον ιατρό.
\item Στην περίπτωση αυτή απαιτούνται τεχνητά δάκρυα, προστατευτικά γυαλιά και οφθαλμολογική εκτίμηση.
\item Αναφέρετε άμεσα ερυθρότητα, πόνο ή θόλωση όρασης στον δεξιό οφθαλμό.
    \end{itemize}
  \item \textbf{Σημείο παρακέντησης (δεξιά παρειά)}
       \begin{itemize}
\item Το σημείο εισόδου της βελόνας παραμένει καλυμμένο με μικρό αυτοκόλλητο για 24--48 ώρες.
\item Ήπιο οίδημα (πρήξιμο) ή εκχύμωση (μελανιά) στην παρειά είναι \textbf{αναμενόμενα} και υποχωρούν σε 5--10 ημέρες.
\item \textbf{Τοπική εφαρμογή παγοκύστης} (τυλιγμένης σε πετσέτα) για 15 λεπτά, 3--4 φορές την ημέρα, τις πρώτες 48 ώρες.
\item Διατηρείτε την περιοχή καθαρή και στεγνή. Πλύσιμο προσώπου με ήπιο σαπούνι επιτρέπεται μετά τις 24 ώρες.
\item Μην πιέζετε και μην μαλάσσετε την περιοχή.
    \end{itemize}
  \item \textbf{Πόνος -- Κεφαλαλγία}
      \begin{itemize}
\item Ήπιος πόνος στη γνάθο, στο αυτί ή στον κρόταφο δεξιά είναι αναμενόμενος τις πρώτες ημέρες.
\item Λαμβάνετε παρακεταμόλη όπως έχει συνταγογραφηθεί.
\item \textbf{Αναφέρετε άμεσα} έντονη κεφαλαλγία, ιδίως εάν συνοδεύεται από πυρετό ή αυχενική δυσκαμψία.
    \end{itemize}
  \item \textbf{Δραστηριότητα -- Ανάπαυση}
      \begin{itemize}
\item Ανάπαυση την πρώτη ημέρα, με σταδιακή επιστροφή στις συνήθεις δραστηριότητες.
\item Αποφύγετε την έντονη σωματική προσπάθεια και την άρση βαρών για 3--5 ημέρες.
\item Αποφύγετε το απότομο σκύψιμο με το κεφάλι χαμηλά τις πρώτες 48 ώρες.
\item Η ήπια βάδιση, με συνοδό, είναι επιθυμητή και συνιστάται.
    \end{itemize}
  \item \textbf{Οδήγηση}

\textbf{Μην οδηγείτε} μέχρι νεότερη ιατρική οδηγία. Τα φάρμακα που λαμβάνετε (Lyrica, Tegretol) επηρεάζουν την εγρήγορση, την όραση και τον χρόνο αντίδρασης.
\end{enumerate}

\newpage
\section*{Εξειδικευμένες Οδηγίες}
\begin{enumerate}[label=\arabic*.]
  \item \textbf{Φαρμακευτική αγωγή για τη νευραλγία}
      \begin{itemize}
\item \textbf{Συνεχίστε κανονικά} τα φάρμακα (Lyrica, Tegretol) στις δόσεις που αναγράφονται στον πίνακα, \textbf{ακόμη και αν ο πόνος έχει υποχωρήσει πλήρως}.
\item \textbf{Μην διακόψετε και μην μειώσετε μόνος/μόνη σας τις δόσεις.} Η απότομη διακοπή, ιδίως του Tegretol, είναι επικίνδυνη.
\item Η \textbf{σταδιακή μείωση} θα ξεκινήσει κατόπιν ιατρικής οδηγίας στην επανεξέταση, εφόσον το αποτέλεσμα είναι σταθερό.
\item Λαμβάνετε τα φάρμακα \textbf{την ίδια ώρα κάθε ημέρα}, με τροφή.
\item Εάν παραλείψετε δόση, μην διπλασιάσετε την επόμενη.
    \end{itemize}
  \item \textbf{Παρενέργειες φαρμάκων -- τι να προσέξετε}
      \begin{itemize}
\item Ζάλη, υπνηλία, θαμπή όραση, διπλωπία, αστάθεια: \emph{συχνά} τις πρώτες ημέρες, υποχωρούν σταδιακά. Βλ.\ οδηγία πρόληψης πτώσεων.
\item Οίδημα στα κάτω άκρα και αύξηση σωματικού βάρους (Lyrica).
\item \textbf{Δερματικό εξάνθημα}: \textbf{διακόψτε το Tegretol και ενημερώστε άμεσα τον ιατρό.}
\item Ναυτία, ανορεξία, σύγχυση, υπνηλία ή γενικευμένη αδυναμία μπορεί να υποδηλώνουν \textbf{χαμηλό νάτριο} από το Tegretol.
\item \textbf{Πυρετός, πονόλαιμος, στοματικά έλκη ή εύκολοι μώλωπες}: ενημερώστε άμεσα τον ιατρό.
\item \textbf{Απαγορεύεται η κατανάλωση αλκοόλ} και ο χυμός γκρέιπφρουτ.
    \end{itemize}
% --- Patient-specific blocks from the source letter (enable if relevant)
%  \item \textbf{Αρτηριακή πίεση}
%      \begin{itemize}
% \item \textbf{Συνεχίστε κανονικά την αντιυπερτασική σας αγωγή} χωρίς αλλαγές.
% \item \textbf{Μετράτε την αρτηριακή πίεση καθημερινά}, την ίδια ώρα, και καταγράφετε τις τιμές. Φέρτε την καταγραφή στην επανεξέταση.
% \item \textbf{Σημαντικό:} το Tegretol μπορεί να \emph{μειώσει τη δράση ορισμένων αντιυπερτασικών φαρμάκων} μέσα σε 1--2 εβδομάδες. Εάν οι τιμές ανέβουν, ενημερώστε τον θεράποντα καρδιολόγο ή παθολόγο σας.
%    \end{itemize}
%  \item \textbf{Νευρολογική παρακολούθηση}
%      \begin{itemize}
% \item Συνεχίστε κανονικά την παρακολούθηση και τη φυσικοθεραπεία για τη νευρολογική σας πάθηση, όπως έχει προγραμματιστεί.
% \item \textbf{Αναφέρετε} οποιαδήποτε νέα ή επιδεινούμενη αδυναμία, σπαστικότητα ή δυσκολία στη βάδιση, ομιλία ή κατάποση.
%    \end{itemize}
  \item \textbf{Ενημέρωση άλλων ιατρών}
      \begin{itemize}
\item Ενημερώστε \textbf{κάθε} ιατρό που σας παρακολουθεί ότι λαμβάνετε Tegretol, λόγω πιθανών αλληλεπιδράσεων.
    \end{itemize}
  \item \textbf{Επούλωση -- Υγιεινή}
        \begin{itemize}
\item Ντους επιτρέπεται μετά τις 24 ώρες, αποφεύγοντας τον έντονο καθαρισμό της περιοχής παρακέντησης.
\item Στοματική υγιεινή με μαλακή οδοντόβουρτσα και ήπιο βούρτσισμα δεξιά.
\item Επισκεφθείτε οδοντίατρο για προληπτικό έλεγχο εντός 1--2 μηνών, ενημερώνοντάς τον για την επέμβαση.
        \end{itemize}
\end{enumerate}

\section*{Επανέλεγχοι}
\begin{description}
\item[Σε 7--10 ημέρες:] κλινική εξέταση, εκτίμηση του αποτελέσματος και της αισθητικότητας του προσώπου.
\item[Σε 1 μήνα:] επανεκτίμηση και \textbf{έναρξη σταδιακής μείωσης της φαρμακευτικής αγωγής}, εφόσον το αποτέλεσμα είναι σταθερό.
\item[Αιματολογικός έλεγχος:] γενική αίματος, ηλεκτρολύτες (νάτριο) και ηπατική βιοχημεία σε \textbf{2--4 εβδομάδες}, λόγω της αγωγής με Tegretol.
\end{description}

\begin{UrgentBox}
\item \textbf{Πυρετό} (Θερμοκρασία $>$ 38\textdegree{}C) με \textbf{έντονη κεφαλαλγία} και \textbf{αυχενική δυσκαμψία} (δυσκολία στην κάμψη του αυχένα προς το στήθος).
\item Φωτοφοβία, ναυτία, έμετοι, σύγχυση ή έντονη υπνηλία.
\item \textbf{Δερματικό εξάνθημα}, ιδίως με φυσαλίδες ή με προσβολή των βλεννογόνων (στόμα, μάτια, γεννητικά όργανα).
\item \textbf{Νέα αδυναμία} στα άκρα ή στο πρόσωπο, διαταραχή ομιλίας, ή απώλεια συνείδησης.
\item Πτώση με τραυματισμό κεφαλής.
\item Ταχέως αυξανόμενο οίδημα στην παρειά ή δυσκολία στο άνοιγμα του στόματος με πυρετό.
\end{UrgentBox}

\textbf{Επικοινωνήστε άμεσα με τον ιατρό εάν εμφανίσετε:}
\begin{itemize}
\item Μούδιασμα στο μέτωπο, στο άνω βλέφαρο ή στη μύτη, ή αίσθηση ότι ο δεξιός οφθαλμός «δεν αισθάνεται».
\item Ερυθρότητα, πόνο, δακρύρροια ή θόλωση όρασης στον δεξιό οφθαλμό.
\item Επίμονο, καυστικό ή ενοχλητικό άλγος σε περιοχή του προσώπου που είναι ταυτόχρονα μουδιασμένη.
\item Επιστροφή του «ηλεκτρικού» πόνου της νευραλγίας.
\item Διπλωπία ή επίμονη θαμπή όραση.
\item Δερματικό εξάνθημα οποιασδήποτε μορφής.
\item Ναυτία, ανορεξία, σύγχυση ή γενικευμένη αδυναμία (πιθανή υπονατριαιμία).
\item Πυρετό, πονόλαιμο, στοματικά έλκη ή ασυνήθιστους μώλωπες.
\item Πληγή ή δήγμα στο εσωτερικό της παρειάς που δεν επουλώνεται.
\item Επίμονη δυσκολία στη μάσηση ή στο άνοιγμα του στόματος.
\item Επιδείνωση της βάδισης, της ισορροπίας ή νέα πτώση.
% \item Σημαντική μεταβολή των τιμών της αρτηριακής πίεσης.
\end{itemize}

\newpage
\begin{center}
\LARGE\textbf{Οδηγίες Φαρμακευτικής Αγωγής}
\end{center}
\vspace{0.5cm}

\begin{MedTable}
\DrugEntry{caps Lyrica}{150mg}{1 x 2}{}{πρωί -- βράδυ, με τροφή. Μην διακόπτεται απότομα}
\DrugEntry{tb Tegretol}{200mg}{1 x 3}{}{με τροφή, σταθερές ώρες. \textbf{Μην διακόπτεται απότομα}}
\DrugEntry{tb Depon}{500mg}{1 ή 2, έως 3 φορές την ημέρα}{}{εφόσον υπάρχει πόνος -- μέγιστο 3g το 24ωρο}
% ---- Drug library: uncomment the lines you need ---------------------
% -- Analgesics
% \DrugEntry{tb Paracetamol}{1000mg}{1 x 3}{10}{}
% \DrugEntry{tb Apotel}{1000mg}{2 ή 3 την ημέρα}{}{παυσίπονο, μέγιστη δόση 3 δισκία το 24ωρο}
% \DrugEntry{tb Depon}{500mg}{1 ή 2, έως 3 φορές την ημέρα}{}{εφόσον υπάρχει πόνος -- μέγιστο 3g το 24ωρο}
% \DrugEntry{tb Tramal}{50mg}{1 x 2 ή 3}{6}{}
% \DrugEntry{tb Zaldiar}{}{2 ή 3 την ημέρα}{}{εφόσον υπάρχει έντονος πόνος}
% -- NSAIDs
% \DrugEntry{tb Movatec}{7.5mg}{1 x 2}{10}{μετά το φαγητό}
% \DrugEntry{tb Movatec}{15mg}{1 x 1}{1}{μετά το φαγητό}
% \DrugEntry{tb Naprosyn}{500mg}{1 x 2}{6}{μετά το φαγητό}
% -- Steroids (tapers)
% \DrugEntry{tb Medrol}{4mg}{1 x 3\newline 1 x 2\newline 1 x 1\newline STOP}{}{για 5 ημέρες\newline για 5 ημέρες\newline για 5 ημέρες}
% \DrugEntry{tb Medrol}{4mg}{2 x 3\newline 1 x 3\newline 1 x 2\newline 1 x 1\newline STOP}{}{για 3 ημέρες\newline για 3 ημέρες\newline για 3 ημέρες\newline για 3 ημέρες}
% \DrugEntry{syr Dexaton}{2mg/5ml}{10cc x 3}{10}{μετά το φαγητό}
% \DrugEntry{tb Dexamethasone}{4mg}{1 x 3}{}{μετά το φαγητό}
% -- Antibiotics
% \DrugEntry{tb Zinadol}{500mg}{1 x 2}{5}{αντιβίωση}
% \DrugEntry{tb Klaricid}{500mg}{1 x 2}{7}{αντιβίωση}
% \DrugEntry{caps Amoxicillin}{500mg}{1 x 3}{7}{μετά το φαγητό}
% -- Gastroprotection
% \DrugEntry{caps Omeprazole}{40mg}{1 x 2}{10}{προ φαγητού}
% \DrugEntry{caps Omeprazole}{20mg}{1 x 2}{10}{15 λεπτά πριν το φαγητό}
% \DrugEntry{caps Mostrelan}{40mg}{1 x 1}{10}{προ φαγητού}
% -- Anticoagulants
% \DrugEntry{inj Clexane}{4000IU}{1 x 1}{10}{αντιπηκτικές ενέσεις υποδορίως, την ίδια ώρα κάθε ημέρα}
% \DrugEntry{inj Innohep}{0,45ml}{1 x 1}{10}{αντιπηκτικές ενέσεις υποδορίως, την ίδια ώρα κάθε ημέρα}
% \DrugEntry{inj Fraxiparine}{2850IU}{1 x 1}{10}{αντιπηκτικές ενέσεις υποδορίως, την ίδια ώρα κάθε ημέρα}
% -- Neuropathic pain, antiepileptics
% \DrugEntry{tb Lyrica}{50mg}{1 x 1 (βράδυ)}{4}{μετά την 4η ημέρα, 1 x 2}
% \DrugEntry{caps Lyrica}{150mg}{1 x 2}{}{πρωί -- βράδυ, με τροφή. Μην διακόπτεται απότομα}
% \DrugEntry{tb Neurontin}{400mg}{2 x 2}{}{}
% \DrugEntry{tb Tegretol}{200mg}{1 x 3}{}{με τροφή, σταθερές ώρες. \textbf{Μην διακόπτεται απότομα}}
% \DrugEntry{caps Cymbalta}{60mg}{1 x 2}{}{}
% \DrugEntry{tb Keppra}{1000mg}{1 x 3}{}{}
% \DrugEntry{tb Depakine}{400mg}{1 x 3}{}{θα γίνει προοδευτική μείωση σε συνεννόηση με τον θεράποντα ιατρό}
% -- Other
% \DrugEntry{sir Duphalac}{5ml}{1 ή 2 την ημέρα}{}{καθαρτικό, ανάλογα με τη λειτουργία του εντέρου -- στόχος: μαλακές κενώσεις χωρίς προσπάθεια}
% \DrugEntry{tb Xanax}{0,25mg}{1 x 2-3}{}{μέχρι 3 την ημέρα επί εκνευρισμού}
% \DrugEntry{tb Escitalopram}{10mg}{1 x 1}{}{}
% \DrugEntry{spray Vioplex-T ή Pulvo47}{}{1 x 2}{}{στο τραύμα, και να καλύπτεται μετά}
% ---------------------------------------------------------------------
\end{MedTable}

% Notes under the table (uncomment the ones you need)
\textbf{Σημείωση 1:} Οι δόσεις των Lyrica και Tegretol \textbf{θα μειωθούν σταδιακά} κατόπιν ιατρικής οδηγίας, εφόσον το αποτέλεσμα της επέμβασης παραμείνει σταθερό. Καμία αλλαγή χωρίς συνεννόηση με τον ιατρό.

\textbf{Σημείωση 2:} Ενημερώστε κάθε ιατρό και φαρμακοποιό ότι λαμβάνετε Tegretol, πριν από τη χορήγηση οποιουδήποτε νέου φαρμάκου.
% \textbf{Σημείωση 3:} Η αντιυπερτασική σας αγωγή συνεχίζεται \textbf{αμετάβλητη}. Καθημερινή μέτρηση και καταγραφή της αρτηριακής πίεσης.

\end{document}
